In this thesis, we studied two type of forecasting problems in genetics that requires modeling approach that is both flexible and can incorporate domain knowledge that are possibly imperfect and incomplete. 

In \cref{ch:doubletrouble}, we developed a Bayesian nonparametric approach in forecasting new and shared variant counts. The approach avoided modeling possibly finite upper bound while allowed the model to grow as more data came in. The method also has built in power-law-like growth as a prior while not requiring users to specify a fixed, known power-law rate but rather learn it from data using an empirical Bayes approach. In \cref{ch:sbirr}, we extended existing Schrödinger bridge approach to incorporate more flexible form of domain knowledge rather than requiring domain scientists to fully specify a model as reference or prior. In \cref{ch:snapmmd}, we developed a least square approach for fitting SDE models that allowed model-based forecasting. In applications to e.g., immune cell activation, our method allowed models with neural network as components while still have bio-chemistry informed structure, allowing robust forecasting. 


A central methodological theme that emerges from this work is the need to balance flexibility and structure in probabilistic modeling. Flexible models—such as nonparametric or neural approaches—are essential for capturing the complexity of modern scientific data. However, in settings with partial observability and limited signal, flexibility alone is insufficient: models that fit the observed data well may fail to generalize or extrapolate in meaningful ways. Conversely, overly restrictive parametric assumptions may impose incorrect structure and lead to biased or misleading conclusions. The methods developed in this thesis aim to navigate this trade-off by combining flexible modeling components with problem-specific structural assumptions. In the variant discovery setting, this takes the form of Bayesian nonparametric models that capture shared and population-specific variation. In the trajectory inference setting, it involves the use of reference dynamics and iterative projection schemes to constrain the space of plausible trajectories. In the forecasting setting, it requires incorporating structure into SDEs to enable stable and interpretable predictions.

Despite being successful, several challenges still remain that should be addressed in future research. In the variant discovery application, although it is mathematically straightforward to generalize the method to more than two populations, the private-shared structure can become combinatorial --- in two population, there are only three kind of variants, namely variants private to population 1, private to population 2 and variants that are shared. In a general $p$ population case there will be $2^p-1$ types. We suspect this brings both statistical and computational challenges. Despite that we are able to avoid \textit{infinite} dimensional integral by having the conjugate structure, we still have to solve a numerical integral problem that grow as number of populations and this can be challenging in more dimensions. Last but not least, we only established power-law bounds rather than exact power-law growth. Future theoretical studies is needed for further understand this class of method theoretically. 

In the trajectory inference and forecasting case, both of our methods in \cref{ch:sbirr} and \cref{ch:snapmmd} require numerically solve SDEs that are computationally intense. In literature there are methods avoiding simulations like \citet{tong2023simulation, zhang2025learning}. However these methods rely on specific reference families rather than allowing general prior knowledge being used. The idea used in \citet{bartosh2025sde} of implicitly parameterize an SDE using marginals and diffusion term can be a promising approach to avoid simulation as the implicitly parameterized model can directly predict snapshot observations. Another direction for development is to back track, i.e., move time backward. Back tracking in SDEs is not the same as forecasting with an inverted time because diffusion processes is not time symmetric. Ideas in \citet{bartosh2025sde} might also be used by define the SDE with time horizons before first observation. Theoretically, when an SDE can be identified using snapshot observation only is still an open problem with some progress in gradient field \citep{lavenant2024toward} and Ornstein-Uhlenbeck processes \citep{guan2024identifying,zhang2025learning} but not in general. Future theoretical research is needed to understand the fundamental limit of this kind of data. 

More broadly, the ideas developed in this thesis are applicable beyond the specific domains considered here. Many problems in science and engineering involve reasoning about partially observed stochastic systems, where the available data provide only indirect information about the underlying processes while useful domain knowledge exists. The perspective of developing methods that are balanced between flexibility and structural provides a general framework for approaching such problems. 

In summary, this thesis argues that reliable inference in many genomics applications requires going beyond traditional approaches that focus solely on fitting observed data. Instead, it requires explicitly modeling the latent structures and dynamics that give rise to the data, while balancing flexibility and structure to enable meaningful extrapolation. By developing methods for forecasting unseen structure, reconstructing latent dynamics, and predicting future behavior, this work contributes to a broader understanding of how to perform statistical inference in settings where data are limited, indirect, and inherently incomplete.
